\documentclass[11pt]{amsart}

\usepackage[T1]{fontenc}
\usepackage{lmodern}
\usepackage[nopatch=footnote]{microtype}
\usepackage[margin=1.1in]{geometry}
\usepackage{amsmath,amssymb,amsthm,mathtools}
\usepackage{booktabs}
\usepackage{xcolor}
\usepackage{hyperref}

\definecolor{linkblue}{HTML}{174A7E}
\hypersetup{
  colorlinks=true,
  linkcolor=linkblue,
  citecolor=linkblue,
  urlcolor=linkblue,
  pdftitle={Bombieri-Vinogradov in Lean},
  pdfauthor={Jonas Whidden}
}

\newtheorem{theorem}{Theorem}[section]
\newtheorem{lemma}[theorem]{Lemma}
\newtheorem{proposition}[theorem]{Proposition}
\newtheorem{corollary}[theorem]{Corollary}
\theoremstyle{definition}
\newtheorem{definition}[theorem]{Definition}
\theoremstyle{remark}
\newtheorem{remark}[theorem]{Remark}

\title{Bombieri-Vinogradov in Lean}
\author{Jonas Whidden}
\date{\today}

\begin{document}

\begin{abstract}
A source-faithful Lean 4 formalization of the Bombieri-Vinogradov theorem in the level-of-distribution form used by modern bounded-prime-gap arguments.
\end{abstract}

\maketitle

\section{Introduction}

Explain the problem's significance, research audience, source relationship,
and the exact scope of the result. Do not claim the project is complete while
the public Challenge/Solution surface is still the template sentinel.

\section{Statement of the problem}

\textbf{The Bombieri-Vinogradov theorem.} For every real theta < 1/2 and every real A >= 1, there exists a real c > 0 such that for every real x >= 3, the sum over 1 <= q <= floor(x\textasciicircum{}theta) of the maximum over reduced residue classes a modulo q of |pi(x; q, a) - pi(x)/phi(q)| is at most c*x/(log x)\textasciicircum{}A.

Define every nonstandard object and state all quantifiers, hypotheses,
endpoints, and exceptional cases.

\section{Relation to the literature}

Give a precise source and prior-formalization account. Distinguish direct
formalization, adaptation, independent proof, background, and original work.

\section{Proof architecture}

Describe the mathematical reductions and how they compose. Replace this
scaffold with the actual proof, not a development chronology.

\section{Formalization and verification}

Map each headline theorem to its Lean declaration. Describe dependency pins,
the Challenge/Solution trust boundary, kernel checks, Comparator/NanoDa replay,
and any exact finite certificates or computations.

\section{Limitations and review status}

State known divergences, omitted material, automation use, and the review that
actually occurred. Mechanical checks are not independent expert review.

\section*{Acknowledgements}

Credit mathematical sources, libraries, tools, contributors, and reviewers.

\begin{thebibliography}{9}

\bibitem{PrimarySource}
Replace with the complete primary-source citation.

\end{thebibliography}

\end{document}
